% ---------------------------------------------------------------------
%  Chapter 13 : Grouped Phase J --- the fully-analytic mode-sum
% ---------------------------------------------------------------------
\chapter{Grouped Phase J: the Fully-Analytic Mode-Sum}\label{ch:grouped-phasej}

The previous chapter walked the \emph{per-diagram} time-domain integrator: for
each typed Feynman diagram, Phase~J takes that diagram's product of retarded
propagators and vertex coefficients and integrates over the internal vertex
times to produce a number $C(t_1, \dots, t_k)$. This chapter is about an
\emph{optimisation} of that loop --- one that is mathematically exact, not an
approximation --- together with the analytic machinery (``Stage~4a'') that
replaces adaptive numerical quadrature with closed-form mode-sums and so lifts
the precision floor from scipy's $\sim\!10^{-8}$ all the way to machine epsilon.

The motivation is brute arithmetic. The enumerator does not emit one diagram per
topology; it emits one diagram per \emph{typed} topology, and a single bare
topology fans out into hundreds of typed copies that differ only in bookkeeping
the integrator can factor out. The grouped path recognises that redundancy, does
the expensive per-topology setup \emph{once}, and sums the integrands of the
whole family before integrating. The result is bit-for-bit the same numbers the
per-diagram path produces (a contract the regression suite enforces to
round-off), produced far faster and, on the analytic path, far more precisely.

Two source files carry the whole story, and we will quote both with
\code{file:line} throughout:
\begin{itemize}
  \item \file{pipeline/\_grouped\_phase\_j.py} --- the \emph{grouping wrapper}
        \code{compute\_correction\_td\_grouped}, which buckets typed diagrams
        into groups and dispatches each group, with a three-tier fallback to the
        per-diagram path when grouping cannot be trusted.
  \item \file{msrjd/integration/time\_domain/grouped\_integral.py} --- the
        \emph{grouped evaluator} \code{integrate\_grouped\_diagram}, which is
        where the sum-then-integrate and the merged-residue mode-sum actually
        happen.
\end{itemize}

\begin{note}
This is a \term{prototype}, opt-in by a single flag. It is \emph{not} wired into
the engine by default: \code{compute\_cumulants} takes a keyword
\code{use\_grouped\_phase\_j} that defaults to \code{False}
(\file{pipeline/compute.py:129}). Pass \code{use\_grouped\_phase\_j=True} to
route Phase~J through the grouped path. Everything downstream --- the $\tau$-grid
evaluation, the plotting, the saving --- is untouched, because the wrapper
returns the \emph{same dict shape} as the per-diagram path. The two paths are
interchangeable behind that one flag.
\end{note}

\section{Why group? The per-diagram redundancy}

To see the redundancy clearly we need two definitions that recur throughout the
chapter. They are the same objects the enumeration and type-assignment chapters
introduced, viewed now from the integrator's side.

\begin{defn}
A \term{prediagram} is the bare graph topology of a Feynman diagram: its
vertices, its \emph{leaves} (the external legs), its edges, and which internal
vertices are integration variables. It carries \emph{no} field-component
information on its edges.

A \term{typed diagram} is a prediagram with two further things fixed: per edge,
\emph{which} propagator matrix entry $(r_i, p_i)$ flows along that edge (i.e.\
which response/physical field-component pair), and the per-vertex coefficients.
\end{defn}

Because a typed diagram only \emph{adds} edge indices and coefficients to a
prediagram, all typed diagrams sharing one prediagram have \emph{identical graph
topology, identical leaves, identical internal vertices, and identical
integration variables}. They differ only in those per-edge $(r_i, p_i)$ indices
and per-vertex coefficients. The module docstring states exactly this
(\file{grouped\_integral.py:9--13}):

\begin{lstlisting}[style=console]
Typed diagrams sharing the same prediagram have identical graph
topology, leaves, internal vertices, and integration variables.
They differ only in per-edge propagator field-type indices (ri, pi)
and per-vertex coefficients.
\end{lstlisting}

The per-diagram integrator does not exploit this. It pays the \emph{entire}
per-diagram setup cost --- the graph walk, the allocation of vertex-time symbols,
the $2^{|E|}$ $\delta$-edge subset bookkeeping (more on that below), the
\code{fast\_callable} compile, the polytope and pole extraction --- once for
\emph{every} typed diagram. The docstring records the measurement that motivated
the whole effort (\file{grouped\_integral.py:18--21}):

\begin{lstlisting}[style=console]
For (k=2, ell=1) multipop the user saw 7736 typed diagrams across
9 prediagrams - ~860 diagrams per prediagram, all paying full setup
cost.
\end{lstlisting}

Eight hundred and sixty typed diagrams per prediagram, each repeating the same
graph walk and the same JIT compile. The setup is identical across all of them;
only a handful of complex numbers differ. That is the redundancy grouping
removes.

\section{Linearity of integration: why grouping is exact}

The mathematical justification is a single line, and it is worth dwelling on
because it is what makes grouping a \emph{drop-in} rather than a tradeoff. The
contribution of one typed diagram is an integral over the internal vertex times;
the contribution of the whole group is the sum of those integrals. Because the
integration \emph{domain} is identical across every typed diagram of one
prediagram --- the domain comes from the shared graph topology, which is the same
for all of them --- you may sum the integrands first and integrate once
(\file{grouped\_integral.py:33--41}):
\begin{equation}
  \sum_{\Gamma}\ \int \dd s_1 \cdots \dd s_m\ \mathcal{I}_\Gamma
  \;=\;
  \int \dd s_1 \cdots \dd s_m\ \sum_{\Gamma}\ \mathcal{I}_\Gamma ,
  \label{eq:linearity}
\end{equation}
where $\mathcal{I}_\Gamma$ is the time-domain integrand of typed diagram
$\Gamma$ and the $s_i$ are the surviving internal-vertex times. This is the
linearity of the integral: it is \emph{exact}, holds term by term, and is not an
approximation in any limit.

So the grouped path builds the prediagram-level scaffolding once, forms the
right-hand-side sum $\sum_\Gamma \mathcal{I}_\Gamma$, and runs a single
integration. The expensive parts (graph walk, symbol allocation, the
$2^{|E|}$ structure, the compile, the polytope) are shared; only the cheap parts
(reading off each diagram's edge residues and coefficients) repeat per diagram.

\begin{note}
\eqref{eq:linearity} is the entire conceptual content of the chapter. Everything
that follows --- the subset expansion, the polytope, the mode-sum, the merged
residue --- is just \emph{how} the right-hand side is built and integrated
efficiently. Keep the picture ``sum the integrands of a family that share a
domain, then integrate once'' in mind and the code reads straightforwardly.
\end{note}

\section{From retarded propagators to a polytope integral}

Before the grouping can buy anything, we need the shape of the integral it is
grouping. This section rebuilds it from the field theory down to the convex
polytope the closed forms integrate over. The construction is shared with the
per-diagram path; the grouped path differs only in the very last step.

\subsection{The retarded propagator in the time domain}

Phase~I (the frequency-domain assembly of the previous part) hands Phase~J a
rational propagator matrix $\hat G(\omega)$ already decomposed into poles and
residues. Concretely it passes a \code{propagator\_data} dict whose keys are:
\begin{description}
  \item[\code{pole\_vals}] the poles $p_k$ of $\hat G(\omega)$. The causality
        filter guarantees $\operatorname{Im}(p_k) > 0$ for every one of them.
  \item[\code{C\_mats}] one residue matrix $C_k$ per pole, so that
        $C_k[i,j]$ is the residue of $\hat G[i,j](\omega)$ at the pole $p_k$.
  \item[\code{D\_delta}] (optional) the $\omega \to \infty$ limits of
        $\hat G$ --- the instantaneous response weights.
\end{description}
The helper \code{build\_G\_t\_matrix} (\file{propagator\_td.py:135}) turns this
into the \term{time-domain retarded propagator}. Its docstring states the
decomposition (\file{propagator\_td.py:140--141}):
\begin{equation}
  \Gret[i,j](t)
  \;=\;
  \underbrace{\text{delta\_coeffs}[i,j]\cdot \delta(t)}_{\text{instantaneous}}
  \;+\;
  \underbrace{\Theta(t)\cdot\Bigl(\sum_k C_k[i,j]\,\ee^{\ii\,p_k\,t}\Bigr)}_{\text{smooth}} .
  \label{eq:GR}
\end{equation}
The first term is the \emph{instantaneous} $\delta$-response (it appears for any
entry whose frequency-domain propagator has a nonzero $\omega\to\infty$ limit,
e.g.\ a response source producing an immediate response at the same time); the
second is the \emph{smooth} pole--residue sum. With the pipeline-wide Fourier
convention, each summand $C_k\,\ee^{\ii p_k t}$ \emph{decays} for $t>0$ and grows
for $t<0$, and the Heaviside $\Theta(t)$ is exactly what makes the retarded
propagator well-defined on the real line (\file{propagator\_td.py:157--161}).

It is convenient to fold the $\ii$ into the rate. Define the \term{mode}
convention the analytic code uses throughout:
\begin{equation}
  \lambda_k \;=\; \ii\,p_k ,
  \qquad
  C_\alpha \;=\; C_k[p_i, r_i] ,
  \label{eq:mode}
\end{equation}
so $\operatorname{Re}\lambda_k < 0$ for a decaying retarded mode, and the smooth
time-domain propagator on an edge reads $\sum_\alpha C_\alpha\,
\ee^{\lambda_\alpha\,\Delta t}$. Here $\Delta t$ is the time difference across the
edge: for an oriented edge $u \to v$, $\Delta t = t_v - t_u$. (The dataclass
\code{EdgeModeSum} records this convention in its own docstring,
\file{final\_integral.py:128--133}, with the index ordering
$G[p_i, r_i] = \langle \varphi_{p_i}\,\tilde n_{r_i}\rangle$ --- physical row,
response column.)

\subsection{One diagram's integral and the $2^{|E|}$ subset expansion}

Give each internal vertex a time --- the integration variables $s_v$ --- and each
leaf its external time, with one leaf (the \term{origin leaf}) pinned to $t=0$.
The contribution of one typed diagram is then
\begin{equation}
  C_\Gamma(t_1,\dots,t_k)
  \;=\;
  (\text{prefactor})\cdot
  \int \prod_v \dd s_v\ \prod_{\text{edges } e} \Gret[p_i, r_i](\Delta t_e) .
  \label{eq:onediagram}
\end{equation}
Each retarded factor is the sum of a $\delta$-part and a $\Theta$-times-smooth
part, \eqref{eq:GR}. Expand the product of $|E|$ such factors over the binary
choice ``$\delta$-piece versus smooth-piece'' on each edge: that gives a sum over
$2^{|E|}$ \term{subsets}. In a given subset, the edges chosen as
``$\delta$-edges'' each carry a $\delta(\Delta t_e)$, which \emph{collapses one
integration variable} via the linear equality $\Delta t_e = 0$; the remaining
``smooth-edges'' each carry $\sum_\alpha C_\alpha \ee^{\lambda_\alpha \Delta t_e}$
together with the retardation constraint $\Delta t_e > 0$ coming from the
Heaviside.

After $\delta$-elimination, suppose $m$ integration variables $s_1,\dots,s_m$
survive. The smooth retardation constraints $\Delta t_e > 0$ are \emph{linear
inequalities} in the $s_i$ and the free external times --- they cut out a
\term{convex polytope} $P$. The subset's contribution is
\begin{equation}
  I_{\text{subset}}
  \;=\;
  (\text{prefactor})\cdot
  \int_P \dd s_1 \cdots \dd s_m\
  \prod_{e\,\in\,\text{smooth}}
  \Bigl[\,\sum_\alpha C^{(e)}_\alpha\, \ee^{\lambda_\alpha \Delta t_e}\,\Bigr] .
  \label{eq:subset}
\end{equation}

\begin{defn}
The integer $m$ --- the number of integration variables surviving
$\delta$-elimination in a subset --- is the single most important quantity in the
evaluator. It selects which closed form is used: $m=0$ (no integral, just a
constraint check), $m=1$ (a 1-D interval), $m=2$ (a polygon), or $m\ge 3$ (a
general polytope). We return to each case in \S\ref{sec:closedforms}.
\end{defn}

\subsection{Pole expansion into mode-sums}

Expand the product of mode-sums in \eqref{eq:subset} into a sum over \term{pole
tuples} $\alpha = (\alpha_1, \dots, \alpha_{n_{\text{smooth}}})$ --- one pole
index per smooth edge. Each tuple contributes a single exponential whose exponent
is \emph{linear} in the $s_i$:
\begin{equation}
  \prod_e \Bigl[\sum_\alpha C_\alpha \ee^{\lambda_\alpha \Delta t_e}\Bigr]
  \;=\;
  \sum_{\bm\alpha}\ \Bigl(\prod_e C_{\alpha_e}\Bigr)\,
  \exp\!\Bigl(\sum_e \lambda_{\alpha_e}\,\Delta t_e\Bigr) .
  \label{eq:poletuple}
\end{equation}
Writing each edge's time difference as a linear form
$\Delta t_e = c_{0,e} + \sum_i a^{\text{int}}_{e,i}\,s_i + \sum_j a^{\text{ext}}_{e,j}\,t^{\text{free}}_j$,
the exponent of \eqref{eq:poletuple} becomes $\gamma_{\bm\alpha} + \sum_i
\alpha^{(i)}_s\,s_i$ with
\begin{equation}
  \alpha^{(i)}_s \;=\; \sum_e \lambda_{\alpha_e}\,a^{\text{int}}_{e,i},
  \qquad
  \gamma_{\bm\alpha} \;=\; \sum_e \lambda_{\alpha_e}\Bigl(c_{0,e} + \sum_j a^{\text{ext}}_{e,j}\,t^{\text{free}}_j\Bigr).
  \label{eq:slopes}
\end{equation}
So the subset contribution becomes a sum over pole tuples, each a constant times
a geometric integral of an exponential-of-linear over the polytope:
\begin{equation}
  I_{\text{subset}}
  \;=\;
  \sum_{\bm\alpha}\Bigl(\prod_e C_{\alpha_e}\Bigr)\,\ee^{\gamma_{\bm\alpha}}
  \int_P \exp\!\Bigl(\sum_i \alpha^{(i)}_s\,s_i\Bigr)\,\dd s .
  \label{eq:subsetfinal}
\end{equation}
The integral $\int_P \exp(\text{linear})\,\dd s$ has a closed form that depends
only on $m$ and the polytope geometry --- it does \emph{not} know which pole tuple
or which diagram produced it. That separation is what the grouped path exploits.

\section{The merged-residue tensor \texorpdfstring{$B_{\bm\alpha}$}{B-alpha}}

Now apply linearity, \eqref{eq:linearity}, to \eqref{eq:subsetfinal}. Take one
prediagram with $N$ typed diagrams indexed by $j$. All $N$ share the same pole
spectrum $\{\lambda_\alpha\}$ (they share \code{propagator\_data}) and the same
polytope (the constraints come from the shared $\Delta t_e$ symbols). Only the
residues $C_\alpha$ and the prefactor $\mathrm{cp}_j$ differ between typed
diagrams --- through their per-edge $(r_i, p_i)$ indices. Summing
\eqref{eq:subsetfinal} over $j$:
\begin{equation}
  I_{\text{subset}}^{\text{group}}
  \;=\;
  \sum_{\bm\alpha}
  \underbrace{\Bigl[\sum_j \mathrm{cp}_j
    \Bigl(\textstyle\prod_{\delta\text{-edge}} \delta_{\mathrm{coeff}\,j,e}\Bigr)
    \prod_{\text{smooth }e} C^{(j)}_{\alpha_e,e}\Bigr]}_{\displaystyle B_{\bm\alpha}}
  \ee^{\gamma_{\bm\alpha}}
  \int_P \exp\!\Bigl(\sum_i \alpha^{(i)}_s s_i\Bigr)\dd s .
  \label{eq:merged}
\end{equation}
Crucially, the $\ee^{\gamma_{\bm\alpha}} \int_P(\cdots)$ factor does \emph{not}
depend on $j$. So the code pre-sums the residue weights over $j$ \emph{inside}
the pole tuple, defining the \term{merged-residue tensor}
(\file{grouped\_integral.py:71--75}):
\begin{equation}
  B_{\bm\alpha}
  \;=\;
  \sum_j\ \Bigl(\mathrm{cp}_j \cdot
    \textstyle\prod_{\delta\text{-edge}} \delta_{\mathrm{coeff}\,j,e}\Bigr)
  \cdot \prod_{\text{smooth }e} C^{(j)}_{\alpha_e,e} .
  \label{eq:Bdef}
\end{equation}

\begin{defn}
The merged-residue tensor $B_{\bm\alpha}$ is the grouped path's per-tuple
coefficient. For the \emph{per-diagram} path the per-tuple coefficient is the
Cartesian product $\prod_e C_{\alpha_e}$ produced by \code{\_enumerate\_pole\_tuples}
(\file{final\_integral.py:529}). The grouped path replaces that single-diagram
product with the sum $B_{\bm\alpha}$ over the whole group. The underlying analytic
integration --- the $\ee^\gamma \int_P$ factor --- is byte-for-byte identical; only
the residue accumulation moved inside the tuple sum
(\file{grouped\_integral.py:79--86}).
\end{defn}

This is the precise sense in which grouping is a drop-in. The per-diagram analytic
evaluators already accept the per-tuple coefficient through an iterator; the
grouped path feeds them the merged $B_{\bm\alpha}$ instead of the Cartesian
product, and the same closed form runs unchanged.

\begin{gotcha}
There is one subtlety, the subject of an open precision investigation, and it
lives entirely in this reordering. The grouped path sums $B_{\bm\alpha}$
\emph{before} the integration, so intra-group cancellations happen at the level
of (possibly large) individual residue products. The per-diagram path sums
\emph{after} integration. By linearity these are equal in exact arithmetic, but
in IEEE float64 they can disagree when individual pole-tuple terms are huge ---
close-paired poles give $O(1/\Delta p)$ residues, products of several such give
$O(10^8)$, and catastrophic cancellation drags them down to an $O(1)$ integral.
The disagreement is documented at \file{final\_integral.py:1690--1700} (the
\code{USE\_POSET\_CAP\_MATCH\_SCIPY} comment block) and affects the $m\ge 3$
case; see \S\ref{sec:precision}.
\end{gotcha}

\section{Closed forms by \texorpdfstring{$m$}{m}}\label{sec:closedforms}

With $B_{\bm\alpha}$ in hand, the only thing left is the geometric integral
$\int_P \exp(\text{linear})\,\dd s$ for each subset. Its closed form is selected
by $m$. The grouped path implements $m=0$ and $m=1$ as \emph{stand-alone twins}
inside \file{grouped\_integral.py}, and routes $m=2$ and $m\ge 3$ through the
shared per-diagram evaluators in \file{final\_integral.py} via a
\code{pole\_tuples} override. All four read the merged residue
$(B_{\bm\alpha}, \text{lambdas})$ pairs the same way.

\subsection{$m=0$: a constraint check}

No integration variable survived $\delta$-elimination, so there is nothing to
integrate --- only the constraints to check and the exponentials to sum. The
helper \code{\_evaluate\_grouped\_m0\_modesum} (\file{grouped\_integral.py:138})
does exactly that:
\begin{lstlisting}
def _evaluate_grouped_m0_modesum(pole_tuples, subset_constraint_data, free_ext_vals):
    import cmath
    for (a_int, a_ext, c0) in subset_constraint_data:
        c_eff = float(c0) + sum(float(a_ext[i]) * float(free_ext_vals[i])
                                for i in range(len(a_ext)))
        if c_eff <= 0:
            return 0.0 + 0.0j
    ...
    total = 0.0 + 0.0j
    for (B_alpha, lambdas) in pole_tuples:
        gamma = 0.0 + 0.0j
        for e in range(len(lambdas)):
            gamma += lambdas[e] * dt_per_edge[e]
        if abs(gamma.real) > _GROUPED_EXP_REAL_LIMIT:
            return None
        try:
            total += B_alpha * cmath.exp(gamma)
        except (OverflowError, ValueError):
            return None
    return total
\end{lstlisting}
It returns $\sum_{\bm\alpha} B_{\bm\alpha}\,\exp(\sum_e \lambda_{\alpha_e}\Delta t_e)$
as a Python \code{complex}; it returns $0$ if any constraint $\Delta t_e > 0$ is
violated (the test \code{c\_eff <= 0} at \file{grouped\_integral.py:155}, the
$\Theta(0)=0$ convention --- the boundary is \emph{outside}); and it returns
\code{None} on overflow, which signals the caller to fall back to numerical
quadrature.

\subsection{$m=1$: an interval, with $\pm\infty$ handled exactly}

One integration variable $s$ survives, on an interval $[L,U]$. The closed form is
elementary,
\begin{equation}
  \int_L^U \ee^{\alpha_s s + \gamma}\,\dd s
  \;=\;
  \frac{\ee^{\alpha_s U + \gamma} - \ee^{\alpha_s L + \gamma}}{\alpha_s},
  \label{eq:m1}
\end{equation}
with the limit $\alpha_s \to 0$ (a constant integrand, value $(U-L)\,\ee^\gamma$)
handled separately. The interesting part is the unbounded endpoints.
\code{\_integrate\_grouped\_m1\_modesum} (\file{grouped\_integral.py:179}) first
resolves the feasible interval from the linear constraints, tracking $\pm\infty$
exactly:
\begin{lstlisting}
    L = -math.inf
    U = +math.inf
    for (a_int, a_ext, c0) in subset_constraint_data:
        a = float(a_int[0]) if a_int else 0.0
        c_eff = float(c0) + sum(float(a_ext[i]) * float(free_ext_vals[i])
                                for i in range(len(a_ext)))
        if abs(a) < 1e-15:
            if c_eff <= 0:
                return 0.0 + 0.0j
            continue
        bound = -c_eff / a
        if a > 0:
            if bound > L: L = bound
        else:
            if bound < U: U = bound
    if L >= U:
        return 0.0 + 0.0j
\end{lstlisting}
A constraint with $|a| < 10^{-15}$ is $s$-independent: if it is violated the
subset is infeasible (return $0$), otherwise it is skipped. A constraint with
$a>0$ tightens the lower bound, $a<0$ the upper bound. When a side is left
unbounded, the corresponding boundary term in \eqref{eq:m1} is \emph{zeroed} only
if the integrand decays in that direction (the sign of
$\operatorname{Re}\alpha_s$ matches); if it would instead diverge, the function
returns \code{None} to force the scipy fallback, which handles the unbounded
endpoint with an adaptive substitution.

\begin{gotcha}
Why $\pm\infty$ rather than clipping at some large finite cap? The docstring
(\file{grouped\_integral.py:196--202}) is explicit, and it is load-bearing for
precision: clipping the interval at $\pm$\code{bbox\_cap} would leave a residual
$\ee^{\alpha_s\cdot\text{bbox\_cap}} \approx 10^{-9}$ boundary term for typical
Hawkes timescales --- measured on the \code{quad\_exp\_k2\_ell0} fixture before
this fix. Treating the genuine $\pm\infty$ exactly removes that error entirely.
The \code{bbox\_cap} argument is still passed in, but on this path it is used only
to \emph{not} clip.
\end{gotcha}

\subsection{$m=2$: a polygon via fan-triangulation}

For two surviving variables the polytope is a convex polygon. The shared
evaluator \code{\_integrate\_2d\_polygon\_modesum} (\file{final\_integral.py:2175})
clips the half-planes, fan-triangulates the polygon from vertex~$0$, and
integrates $\exp(\alpha x + \beta y)$ over each triangle using the unit-triangle
closed form $J(p,q)$ (\file{final\_integral.py:356}):
\begin{equation}
  J(p,q) \;=\; \int_0^1\!\dd u \int_0^{1-u}\!\dd w\ \ee^{p\,u + q\,w} .
  \label{eq:J}
\end{equation}
Its docstring gives the derivation (do the $w$-integral first, then split and
integrate in $u$) and notes a 4th-order Taylor fallback when any of $|p|$, $|q|$,
$|p-q|$ is tiny --- the same catastrophic-cancellation guard that recurs across
the analytic code. The grouped path calls this evaluator with the merged residues
supplied as \code{pole\_tuples=grouped\_pole\_tuples}, a dummy
\code{smooth\_edge\_modes=list(dummy\_modes)} (explained in
\S\ref{sec:dummymodes}), and \code{prefactor\_complex=1.0+0.0j} --- because the
prefactor is already folded into $B_{\bm\alpha}$.

\subsection{$m\ge 3$: a causal-poset chain simplex}

For three or more variables the polytope is handled by a \term{causal-poset
chain-simplex} decomposition in \code{\_integrate\_nd\_polytope\_poset\_modesum}
(\file{final\_integral.py:1726}). The retardation constraints define a partial
order $s_u < s_v$ on the vertex times; each linear extension of that poset is a
nested simplex $L \le s_{\sigma(0)} \le \cdots \le s_{\sigma(m-1)} \le U$, and the
integrator sums $\int \exp(\sum \lambda s)$ over every nested simplex. The grouped
path uses it with the same \code{pole\_tuples} override. This is the path with the
known float64 cancellation issue (\S\ref{sec:precision}).

\begin{note}
The two flags \code{USE\_POLYGON\_M2\_INTEGRATOR}
(\file{final\_integral.py:287}) and \code{USE\_POSET\_INTEGRATOR}
(\file{final\_integral.py:1661}) gate the $m=2$ and $m\ge 3$ analytic paths
respectively. The grouped file imports the \emph{module}, not the flag values
(\code{import \ldots final\_integral as \_fi\_mod}, \file{grouped\_integral.py:119}),
so each per-call read picks up a notebook override live. Flip
\code{\_fi\_mod.USE\_POSET\_INTEGRATOR = False} in a notebook and the grouped
$m\ge 3$ analytic path silently turns off and routes to scipy.
\end{note}

\section{The grouping wrapper}

We now walk the two files top to bottom, starting with the wrapper in
\file{pipeline/\_grouped\_phase\_j.py}. Its job is to bucket typed diagrams into
groups, dispatch each group to the evaluator, and --- the part that earns its
keep --- fall back to the per-diagram path whenever the grouped result cannot be
trusted.

\subsection{The skip-status guard}

Before any logic, the module declares the set of subset-level statuses that mean
``the grouped evaluator silently \emph{dropped} this subset rather than
integrating it'' (\file{\_grouped\_phase\_j.py:49}):
\begin{lstlisting}
_SUBSET_SKIP_STATUSES = frozenset({
    'shotnoise_skipped_in_prototype',
    'dt_sym_mismatch_skipped',
})
\end{lstlisting}
This frozenset is the heart of the wrapper's correctness story. A group can return
an overall \code{status='ok'} and yet have a \emph{subset} diagnostic carrying one
of these codes --- which means a piece of the integrand was dropped and the group
total is numerically \emph{wrong}. The comment (\file{\_grouped\_phase\_j.py:42--48})
calls this ``the production bug we're fixing'' and lists the statuses explicitly
so that a \emph{new} skip status added later cannot silently regress the guard.

\begin{gotcha}
\textbf{If you add a new skip status to the evaluator, you MUST add it here.}
Otherwise the wrapper will accept an incomplete integrand as if it were
complete. The isolation test cited in the comment found
``5/7 multi-diagram groups dropped to 0'' at spike-reset $k=1$, $\ell=2$ ---
exactly the failure mode this guard exists to catch.
\end{gotcha}

\subsection{The entry point and the group key}

The entry point is \code{compute\_correction\_td\_grouped}
(\file{\_grouped\_phase\_j.py:55}):
\begin{lstlisting}
def compute_correction_td_grouped(
    typed_diagrams, prefactors, propagator_data, k,
    num_params=None, ext_time_vars=None,
    origin_leaf_idx=0, external_fields=None,
):
\end{lstlisting}
It takes a flat list of \code{typed\_diagrams} at one loop order, a parallel list
of \code{prefactors} (the \code{combined\_prefactor} per diagram), the Phase~I
\code{propagator\_data}, the leg count \code{k}, optional numeric substitutions
\code{num\_params}, the external time symbols \code{ext\_time\_vars} (auto-built as
$t_1,\dots,t_k$ if \code{None}, \file{\_grouped\_phase\_j.py:72--76}), the pinned
\code{origin\_leaf\_idx}, and the canonical \code{external\_fields} list. It
returns the same dict shape as \code{compute\_correction\_td}.

The grouping itself is a \code{defaultdict} keyed on a pair
(\file{\_grouped\_phase\_j.py:96--100}):
\begin{lstlisting}
groups = defaultdict(lambda: {'tds': [], 'cps': []})
for td, pf in zip(typed_diagrams, prefactors):
    key = (id(td.prediagram), _ext_legs_key(td))
    groups[key]['tds'].append(td)
    groups[key]['cps'].append(pf)
\end{lstlisting}

\begin{defn}
\code{collections.defaultdict} is a \code{dict} subclass that auto-creates a
default value for any missing key, by calling a zero-argument factory you supply.
Here the factory \code{lambda: \{'tds': [], 'cps': []\}} makes a fresh
two-list bucket the first time a key is seen, so
\code{groups[key]['tds'].append(td)} works without a membership check. It is the
idiomatic Python way to bucket items into groups.
\end{defn}

The key has two components. The first, \code{id(td.prediagram)}, is the prediagram
identity --- Python's built-in \code{id()} returns a unique integer per object, and
the enumerator emits the \emph{same} prediagram object for all typed diagrams of one
topology, so identity is a valid grouping key. The second, the nested
\code{\_ext\_legs\_key(td)} (\file{\_grouped\_phase\_j.py:87}), returns a hashable
sorted tuple of the external-legs mapping:
\begin{lstlisting}
def _ext_legs_key(td):
    return tuple(sorted((lf, fld) for lf, fld in td.external_legs.items()))
\end{lstlisting}
Subdividing by external-legs is \emph{required}, not cosmetic. Two typed diagrams
that differ only in which leaf carries which external field have different Wick
mappings and different compensation factors inside the evaluator, so they cannot
share a sum (\file{\_grouped\_phase\_j.py:78--86}). The per-diagram path handles
each independently; the grouped path subdivides at this finer key to stay safe.

\subsection{The three-tier dispatch}

For each group the wrapper tries the grouped evaluator and falls back in tiers
(\file{\_grouped\_phase\_j.py:125--228}):

\begin{description}
  \item[Tier 1 --- accept the grouped result.] Call
        \code{integrate\_grouped\_diagram}. If it returns \code{status='ok'}
        \emph{and} no subset diagnostic is in \code{\_SUBSET\_SKIP\_STATUSES}
        (the \code{has\_skipped\_subset} scan at
        \file{\_grouped\_phase\_j.py:147--153}), accept it: append
        \code{result['contribution']} to the running list of group callables and
        record metadata with \code{'handled\_by': 'grouped\_evaluator'}.
  \item[Tier 2 --- per-diagram fallback on a silent subset skip.] If
        \code{status='ok'} but a subset \emph{was} silently dropped, re-compute
        the \emph{entire group} via the per-diagram path
        (\code{\_perdiag\_fallback}, \file{\_grouped\_phase\_j.py:107}), append
        that callable, and record \code{'handled\_by': 'perdiag\_fallback'} with
        the skip reasons. This trades the grouped speedup on that group for the
        correct numbers --- exactly what \code{use\_grouped\_phase\_j=False} would
        have produced.
  \item[Tier 3 --- per-diagram fallback on a group-level failure.] If the grouped
        evaluator failed outright (e.g.\ early validation), also fall back to
        per-diagram. If \emph{that} throws too, record a \emph{loud}
        \code{'skipped'} entry (\file{\_grouped\_phase\_j.py:213--228}) rather than
        silently dropping the whole group.
\end{description}

The nested \code{\_perdiag\_fallback(tds, cps)} (\file{\_grouped\_phase\_j.py:107})
simply calls the per-diagram \code{compute\_correction\_td} on the group's diagrams
and returns \code{(pd\_result['total\_C'], len(tds))}, isolating the fallback so its
returned callable has the same \code{(*ext\_time\_values) -> complex} signature as
the grouped \code{contribution}.

\subsection{Assembling the result}

Finally the wrapper builds the two callables the downstream code expects
(\file{\_grouped\_phase\_j.py:231--241}):
\begin{lstlisting}
def total_C(*ext_time_values):
    total = 0.0 + 0.0j
    for fn in group_callables:
        total = total + complex(fn(*ext_time_values))
    return total

def total_C_batch(tau_points, parallel=False, n_workers=None):
    return [complex(total_C(*pt)) for pt in tau_points]
\end{lstlisting}
\code{total\_C} sums every group callable; \code{total\_C\_batch} serially
evaluates \code{total\_C} at each $\tau$-point. The returned dict carries
\code{total\_C}, \code{total\_C\_batch}, the grouped bookkeeping (\code{groups},
\code{skipped\_kernel\_ids}, \code{fallback\_to\_perdiag}), and \code{ext\_time\_vars};
it sets \code{eval\_per\_diagram\_batch} to \code{None} and
\code{delta\_contributions} to \code{[]} --- both unsupported in the prototype
(\file{\_grouped\_phase\_j.py:243--252}).

\begin{gotcha}
\code{total\_C\_batch} ignores its \code{parallel} and \code{n\_workers}
arguments --- the prototype is serial. It keeps the arguments only to match the
per-diagram API signature. This is deliberate and safe: as the codebase memory
records, fork-based parallelism crashes a Jupyter kernel on the user's macOS
machine, so serial-only is the correct default here.
\end{gotcha}

\section{The grouped evaluator, walked}

The evaluator \code{integrate\_grouped\_diagram} (\file{grouped\_integral.py:283})
is where \eqref{eq:linearity} and \eqref{eq:merged} are realised. We follow its
control flow.

\subsection{Validation and shared scaffolding}

It first validates the group (\file{grouped\_integral.py:328--358}): an empty list
fails; a length mismatch between diagrams and prefactors fails; and --- the
defining invariant --- \emph{all} typed diagrams must share the same prediagram
object by identity (\file{grouped\_integral.py:349--350}):
\begin{lstlisting}
pd0 = typed_diagrams[0].prediagram
if not all(td.prediagram is pd0 for td in typed_diagrams):
    return {'status': 'failed', ...}
\end{lstlisting}
Then it builds the shared scaffolding once (\file{grouped\_integral.py:360--379}):
the loop number $L = |E| - |V| + 1$ from \code{\_loop\_number\_from\_graph}
(\file{final\_integral.py:2401}); the Sage graph \code{D = td0.prediagram[0]} and
its leaves \code{td0.prediagram[2]}; and the time-domain propagator matrix
\code{G\_t\_obj = build\_G\_t\_matrix(propagator\_data, t\_sym, num\_params)} built
\emph{once} for the whole group.

\subsection{Wick enumeration and compensation}

Next it enumerates the external-Wick mappings and computes the compensation
divisor (\file{grouped\_integral.py:381--452}). If \code{external\_fields} is
provided and matches the leaf count, it buckets canonical external positions and
leaves by field and takes \code{itertools.permutations} within each field,
building the Cartesian product of per-field permutations into \code{\_all\_mappings}.
If any field's counts mismatch it falls back to the identity mapping with
\code{\_compensation = 1}. Otherwise it computes the \emph{exact} compensation via
\code{external\_wick\_compensation} on every group member and asserts they all
agree (\file{grouped\_integral.py:434--449}); a mismatch returns
\code{status='failed'}, because a shared divisor would mis-weight the group. We
unpack the compensation in \S\ref{sec:compensation}.

\subsection{Vertex times and per-diagram edge info}

For the first mapping it assigns each leaf its external time, pinning the origin
leaf to \code{SR(0)}, and allocates a fresh symbol \code{s\_v...} per internal
vertex into \code{integration\_vars} (\file{grouped\_integral.py:454--469}). A
per-diagram scan handles \code{NoiseSourceType} vertices that introduce extra
$\tau$ integration variables (\file{grouped\_integral.py:471--520}); for plain
groups this loop is a no-op and the analytic path stays enabled.

Then comes \emph{the only genuinely per-diagram work}
(\file{grouped\_integral.py:522--591}). For each typed diagram and each edge
$(u,v,\text{lbl})$ of the shared prediagram, it looks up $(r_i, p_i)$ via
\code{\_lookup\_prop\_indices} (\file{final\_integral.py:5230}), resolves the head
and tail times, and builds the per-edge data:
\begin{lstlisting}
dt = SR(t_v - t_u)
delta_c = G_t_delta_coeff(G_t_obj, pi, ri)
smooth_factor = G_t_entry(G_t_obj, pi, ri, dt, include_heaviside=False)
edge_info.append({'u': u, 'v': v, 'lbl': lbl, 'ri': ri, 'pi': pi,
                  'dt_sym': dt, 'delta_coeff': delta_c,
                  'smooth_factor': smooth_factor})
\end{lstlisting}
Note \code{include\_heaviside=False}: causality is enforced by the polytope
constraints, not by an explicit $\Theta$. (\code{G\_t\_entry} reads the smooth
$[p_i, r_i]$ entry --- the transpose of the kernel's $(r_i, p_i)$ convention, which
both Phase~I and Phase~J must agree on.) It then coerces the diagram's
\code{combined\_prefactor} to \code{SR} and substitutes \code{num\_params}.

\subsection{The analytic mode-sum precomputation}

The analytic path is enabled only if \code{USE\_GROUPED\_ANALYTIC\_MODESUM} is set,
\code{pole\_vals} and \code{C\_mats} are present, \emph{and} no diagram has a
non-empty \code{vertex\_leg\_time} (\file{grouped\_integral.py:602--607}):
\begin{lstlisting}
grouped_analytic_enabled = (
    USE_GROUPED_ANALYTIC_MODESUM
    and pole_vals is not None
    and C_mats is not None
    and not any(vlt for vlt in per_td_vertex_leg_time)
)
\end{lstlisting}
When enabled it precomputes, once, the per-pole rates
$\lambda_k = \texttt{complex(CDF(SR(p)))}\cdot 1j$
(\file{grouped\_integral.py:612}) and the per-diagram/per-edge/per-pole residue
tensor \code{residues\_per\_td\_edge} (\file{grouped\_integral.py:619--628}). These
are reused across all $2^{|E|}$ subsets.

\subsection{The subset loop and merged-residue construction}

The core is the loop over \code{subset\_bits} in \code{range(2**n\_edges)}
(\file{grouped\_integral.py:654--1033}). For each subset it: splits edges into
$\delta$-edges (bit set) and smooth-edges (bit clear); drops any diagram with a
near-zero $\delta$-coefficient on a $\delta$-edge; checks that all contributing
diagrams agree on the symbolic $\Delta t$ of each $\delta$-edge (a mismatch
records \code{dt\_sym\_mismatch\_skipped} and skips --- one of the two
\code{\_SUBSET\_SKIP\_STATUSES}); performs $\delta$-elimination with
\code{sage\_solve}; and skips residual shot-noise $\delta$-spikes with status
\code{shotnoise\_skipped\_in\_prototype} (the other guarded status).

It then builds the \emph{summed} subset integrand --- the right-hand side of
\eqref{eq:linearity} --- as a single \code{SR} expression
(\file{grouped\_integral.py:759--770}):
\begin{lstlisting}
subset_factor_group = SR(0)
for td_idx in contributing_td:
    ei = per_td_edge_info[td_idx]
    cp = per_td_cp[td_idx]
    term = cp
    for i in delta_edges:
        term = term * SR(ei[i]['delta_coeff'])
    for i in smooth_edges:
        term = term * ei[i]['smooth_factor']
    term = term.subs(substitutions)
    subset_factor_group = subset_factor_group + term
\end{lstlisting}
After extracting the linear constraint data $(a^{\text{int}}, a^{\text{ext}}, c_0)$
for each smooth edge (\file{grouped\_integral.py:830--847}) and compiling a single
\code{fast\_callable} over the summed integrand (used only on the scipy fallback),
it constructs the merged-residue pole tuples --- the realisation of
\eqref{eq:Bdef} (\file{grouped\_integral.py:925--948}):
\begin{lstlisting}
grouped_pole_tuples = []
idx_alpha = [0] * n_smooth
while True:
    B_alpha = 0.0 + 0.0j
    for j in range(len(merged_pf_per_j)):
        prod = merged_pf_per_j[j]
        for ee in range(n_smooth):
            prod *= smooth_residues_per_j[j][ee][idx_alpha[ee]]
        B_alpha += prod
    lambdas = tuple(lambdas_per_pole[idx_alpha[ee]] for ee in range(n_smooth))
    grouped_pole_tuples.append((B_alpha, lambdas))
    # advance the multi-index alpha ...
\end{lstlisting}
The outer \code{while} enumerates the multi-index $\bm\alpha$ over per-edge poles;
the inner sum over $j$ accumulates $B_{\bm\alpha} = \sum_j (\mathrm{cp}_j
\prod_\delta \delta) \prod_{\text{smooth}} C^{(j)}_{\alpha_e}$, with
\code{merged\_pf\_per\_j} the per-diagram prefactor-times-$\delta$-coefficients
built at \file{grouped\_integral.py:904--913}. Each $(B_{\bm\alpha},
\text{lambdas})$ pair is exactly what the closed-form evaluators consume.

\subsection{The dummy modes}\label{sec:dummymodes}

The per-diagram $m=2$ and $m\ge 3$ evaluators expect a list of \code{EdgeModeSum}
objects of length $n_{\text{smooth}}$ --- a length check. On the analytic path
those modes are never read (the real weights are in \code{pole\_tuples}), so the
grouped code builds throwaway \code{EdgeModeSum} instances purely to satisfy the
check (\file{grouped\_integral.py:949--961}). Their \code{modes} come from the
first diagram's residues.

\begin{gotcha}
Do not mistake \code{grouped\_dummy\_modes} for the actual weights. They carry the
first diagram's residues so the \code{n\_smooth} length check passes, but when
\code{pole\_tuples} is supplied the evaluator reads its coefficients from there,
not from the modes. The real merged weights live in \code{grouped\_pole\_tuples}.
\end{gotcha}

\subsection{The Wick-permutation sum and the divisor}

The final \code{contribution(*ext\_time\_values)} closure
(\file{grouped\_integral.py:1049--1074}) loops over the Wick permutations: for
each it permutes the external times, time-shifts so the permuted origin leaf
returns to $0$ (the integrand is a function of time \emph{differences}, so this
symmetrises asymmetric integrands and is a no-op for symmetric ones), extracts the
free-time values, and sums every subset closure. Finally it divides by the
compensation:
\begin{lstlisting}
        for cfn in subset_contributions:
            total = total + complex(cfn(free_vals))
    return total / _comp
\end{lstlisting}
That \code{/ \_comp} at \file{grouped\_integral.py:1074} is the external-Wick
compensation, the subject of the next section.

\section{External-Wick compensation}\label{sec:compensation}

The integrators enumerate \emph{every} field-respecting assignment of canonical
external positions to diagram leaves and sum the integrand over all of them. Some
of those assignments are redundant --- they map the diagram to itself via a graph
automorphism --- so the raw sum over-counts each \emph{distinct} pinned-external
diagram. The exact divisor that corrects this is the \term{orbit--stabilizer
index}
\begin{equation}
  \mathrm{comp}
  \;=\;
  \frac{|\Aut(\Gamma,\ \text{leaves free})|}{|\Aut(\Gamma,\ \text{leaves fixed})|},
  \label{eq:comp}
\end{equation}
computed by \code{external\_wick\_compensation} (\file{symmetry.py:343}). By
orbit--stabilizer, two assignments yield the same pinned-external diagram iff they
differ by a leaf-permuting automorphism, and the stabilizer of any assignment is
exactly $\Aut(\Gamma, \text{leaves fixed})$; so the mapping sum counts each
distinct class \eqref{eq:comp} times, and dividing by it recovers the sum over
distinct diagrams the Feynman rules require. The grouped path applies this divisor
once, at \file{grouped\_integral.py:1074}.

\begin{note}
\code{external\_wick\_compensation} computes \eqref{eq:comp} via two calls to
\code{\_automorphism\_order}, which builds a coloured incidence digraph and asks
SageMath for the order of its automorphism group. If $|\Aut_{\text{free}}|$ is not
divisible by $|\Aut_{\text{fixed}}|$ --- impossible in exact group theory, since the
fixed group is a subgroup of the free one, but reachable as a Sage edge case --- it
raises a \code{RuntimeError} rather than return a wrong divisor
(\file{symmetry.py:368--381}). Failing loudly is the whole point of a function
that exists to prevent silent mis-weighting.
\end{note}

\subsection{Why the old heuristic broke at \texorpdfstring{$k\ge 3$}{k>=3}}

The history here is instructive, because it explains a function name that still
appears in the codebase. The previous implementation \emph{approximated} the index
\eqref{eq:comp} as a product of factorials, $\prod N_{\text{sig,field}}!$, over
leaves grouped by the \code{vertex\_role\_signature} (\file{symmetry.py:247}) of
their attachment vertex. The \code{vertex\_role\_signature} is a \emph{depth-1}
graph invariant --- a vertex's type, coefficient, bigrade, external-leaf attachment
pattern, and internal-edge incidence pattern.

That heuristic \emph{over-divides} whenever two attachment vertices share a role
signature (a shallow, depth-1 invariant) but are \emph{not} related by any
automorphism. The documented example is the $k=4$ OU$+\varepsilon x^3$ one-loop
cascades, where three noise vertices all read ``noise feeding a cubic'' at depth~1
yet hang off structurally distinct cubics. The resulting deficits broke
\emph{every} $k\ge 3$ one-loop cumulant while leaving $k=2$ exact --- at $k=2$ the
heuristic happens to coincide with the true index. It was replaced by the exact
Sage automorphism ratio (\file{symmetry.py:368--381}), validated against the exact
Boltzmann series $\kappa_4 = -6\varepsilon + 126\varepsilon^2$.

\begin{gotcha}
\code{vertex\_role\_signature} still exists and is still referenced elsewhere
(\file{final\_integral.py:2573}), but it is \emph{not} the current compensation
mechanism and is not on the grouped path's numerical route. Do not reach for it
when reasoning about the divisor; the grouped path uses the exact
\code{external\_wick\_compensation}. The signature survives only as the historical
basis of the superseded heuristic.
\end{gotcha}

\section{The tooling, for the physicist}

This subsystem leans almost entirely on \term{SageMath}, with a fallback into
\code{scipy}. There is no direct use of nauty, sympy, numba, or networkx in the
two primary files --- though Sage embeds several of those internally. Here is each
piece, from the ground up, with the actual import and call sites.

\subsection{SageMath: \texttt{SR}, \texttt{fast\_callable}, \texttt{CDF}, \texttt{solve}}

\term{SageMath} is a large open-source mathematics system built on Python. It
bundles dozens of specialised libraries (a symbolic-algebra engine, a numerical
stack, graph-automorphism engines, number-theory backends) behind one Python API.
When the code writes \code{from sage.all import ...} it is importing Sage's
``everything'' namespace. The grouped file imports four names
(\file{grouped\_integral.py:101}):
\begin{lstlisting}
from sage.all import SR, fast_callable, CDF, solve as sage_solve
\end{lstlisting}

\begin{description}
  \item[\code{SR} --- the Symbolic Ring.] Sage's ring of symbolic expressions
        (its analogue of SymPy's \code{Symbol}/\code{Expr}, but a different
        engine). \code{SR.var('x')} makes a symbolic variable; \code{SR(3)} wraps
        a number; arithmetic builds expression trees you can \code{.subs()},
        \code{.expand()}, \code{.coefficient()}, \code{.variables()}, and so on.
        The grouped path uses it everywhere: to allocate the internal time symbol
        \code{t\_sym = SR.var('\_t\_td\_')} (\file{grouped\_integral.py:378}), the
        per-vertex time symbols (\file{grouped\_integral.py:467}), the symbolic
        per-edge $\Delta t$ via \code{dt = SR(t\_v - t\_u)}
        (\file{grouped\_integral.py:549}), and the summed subset integrand
        accumulated with \code{+} and \code{*}
        (\file{grouped\_integral.py:760--770}).
  \item[\code{fast\_callable} --- a JIT-style expression compiler.] A symbolic
        \code{SR} expression evaluated naively re-walks its expression tree on
        every call, which is slow. \code{fast\_callable(expr, vars=[...],
        domain=CDF)} compiles it into a fast closure taking the listed variables
        as positional arguments. It is Sage's analogue of SymPy's
        \code{lambdify}. The grouped path builds it once per subset over the
        \emph{summed} integrand (\file{grouped\_integral.py:815--818}) --- one
        compile instead of $N$ --- but only the scipy fallback ever calls the
        result; the analytic path never touches it.
  \item[\code{CDF} --- the Complex Double Field.] Sage's complex numbers backed
        by hardware \code{double} precision. \code{complex(CDF(SR(expr)))} is the
        idiom for ``evaluate this symbolic expression to a fast machine-precision
        Python \code{complex}.'' It is the \code{fast\_callable} domain and the
        cast target for poles and residues, e.g.\ \code{complex(\_CDF(SR(p)))*1j}
        for a rate (\file{grouped\_integral.py:612}).
  \item[\code{solve} (as \code{sage\_solve}) --- symbolic equation solver.] Sage's
        \code{solve(eqn, var, solution\_dict=True)} solves an equation
        symbolically and returns a list of solution dicts. The grouped path uses
        it to \emph{eliminate} an integration variable on each $\delta$-edge: the
        $\delta$-edge imposes the linear equation $\Delta t_e = 0$, and solving it
        for one surviving variable removes that variable from the integral
        (\file{grouped\_integral.py:719--722}).
\end{description}

\subsection{nauty / bliss, reached through Sage}

The compensation factor needs the \emph{order} of the automorphism group of a
coloured graph. Sage computes this with its built-in graph-isomorphism engine
(historically nauty/bliss-class algorithms). The grouped path reaches it
indirectly, through \code{external\_wick\_compensation} $\to$
\code{\_automorphism\_order} $\to$ Sage's \code{automorphism\_group().order()} on
a coloured digraph. The novice takeaway: \emph{nauty/bliss answer ``how many ways
can I relabel this graph's vertices and leave it unchanged, respecting the
colours?'', and that count is the divisor \eqref{eq:comp} that prevents the
external-leaf permutation sum from over-counting.}

\subsection{scipy, the numerical fallback}

\code{scipy} is the standard Python scientific-computing library;
\code{scipy.integrate.nquad} performs \term{adaptive multidimensional numerical
quadrature} --- it samples the integrand on an adaptively-refined grid until a
relative tolerance (default $\sim\!10^{-8}$) is met. The grouped path does not
import scipy directly, but its fallback routes through \code{\_integrate\_polytope}
(\file{final\_integral.py:4372}), which is scipy-backed. The grouped path only
reaches scipy when an analytic evaluator returns \code{None} (overflow, an
unbounded-divergent endpoint, a non-rational kernel). This is the slower,
lower-precision route; Stage~4a exists precisely to avoid it.

\begin{note}
What is \emph{not} used here: \code{numba} (the analytic mode-sum is pure
Python-level complex arithmetic, fast enough without it; numba lives in the
spatial path); \code{networkx} (all graph work goes through Sage \code{Graph}
objects carried on the prediagram); and \code{sympy} (Daedalus's CAS is Sage's
\code{SR}, not SymPy). The standard-library helpers are \code{defaultdict} (for
the group buckets), \code{itertools.permutations} (for the Wick mappings), and
\code{cmath}/\code{math} (the complex exponential and exact $\pm\infty$ tracking
in the hot helpers).
\end{note}

\section{Precision: analytic versus scipy, grouped versus per-diagram}\label{sec:precision}

There are two precision stories here, and they should not be confused.

The first is \term{analytic versus scipy}, and the grouped path wins it cleanly.
On the analytic merged-residue path every subset integral is a closed-form
mode-sum, evaluated in machine-precision complex arithmetic, so the precision
floor is machine $\varepsilon$ ($\sim\!10^{-15}$). The scipy fallback is adaptive
quadrature at $\sim\!10^{-8}$ relative. Stage~4a (the master switch
\code{USE\_GROUPED\_ANALYTIC\_MODESUM = True} at \file{grouped\_integral.py:130})
is what moved the grouped path onto the analytic route by default.

The second is \term{grouped versus per-diagram}, and here lies the open issue. By
\eqref{eq:linearity} the two paths are equal in exact arithmetic, and the codebase
treats agreement to floating-point round-off as a \emph{correctness contract}: the
regression suite \file{tests/test\_grouped\_vs\_perdiag.py} computes the sum of $N$
per-diagram results and the single grouped result and asserts they match to
$\sim\!10^{-14}$ relative on the analytic path. The working fixture for this is
\code{quad\_exp\_k2\_ell0}.

\begin{gotcha}
\textbf{Open precision bug at $m\ge 3$ with close-paired poles.} The grouped path
sums $B_{\bm\alpha}$ \emph{before} the chain-simplex integration; the per-diagram
path sums \emph{after}. When close-paired poles produce $O(1/\Delta p)$ residues
and a tuple multiplies several of them into $O(10^8)$, the two summation orders
disagree in float64 even though they agree exactly. The spike-reset $k=2$,
$\ell=1$ case shows a $\sim\!4\times$ discrepancy (per-diagram $+2.66\times10^{-3}$
versus grouped $+6.38\times10^{-4}$). The \code{USE\_POSET\_CAP\_MATCH\_SCIPY}
analysis (\file{final\_integral.py:1690--1700}) concludes the bounding-box cap is
\emph{not} the dominant cause and points to a bookkeeping difference in pole-tuple
construction between the two paths. An attempted fix using extended-precision
accumulation (\code{USE\_POSET\_MPMATH\_ACCUMULATION}, off by default) ``did NOT
address the spike-reset bug.'' The authoritative writeup is
\file{docs/m\_ge3\_precision\_bug\_audit.md}.
\end{gotcha}

\begin{note}
A separate, easily-confused failure is \emph{not} an integrator bug at all. Four
\code{spike\_reset\_*} cases in \file{tests/test\_grouped\_vs\_perdiag.py} fail
because a fixture still passes a \code{taug} parameter that the theory file
dropped when it switched to a delta synaptic kernel; the failure happens in the
pole/residue computation \emph{before} any Phase~J integration runs. That is
fixture/theory drift, not a Phase~J defect. Use \code{quad\_exp\_k2\_ell0} as the
trustworthy regression fixture.
\end{note}

\section{Limitations and the fallback matrix}

The grouped path is a prototype, and it is honest about what it does not handle.
In each of the following cases the grouped evaluator either disables its analytic
route (falling to scipy within the grouped path) or signals the wrapper to fall
back to the per-diagram path entirely. Knowing the matrix tells you when you must
use \code{use\_grouped\_phase\_j=False}.

\begin{description}
  \item[Noise-source kernels disable the analytic path.] If \emph{any} diagram in
        the group has a non-empty \code{vertex\_leg\_time} --- a
        \code{NoiseSourceType} cumulant kernel --- the integrand contains
        non-rational factors and the merged-residue closed form cannot fire. The
        subset falls through to \code{fast\_callable} $+$ scipy
        (\file{grouped\_integral.py:602--607}). The numbers are still correct,
        just at scipy precision.
  \item[Shot-noise $\delta$-spikes are skipped.] A residual external-time equality
        after $\delta$-elimination ($\tau=0$ same-population legs in heterogeneous
        Hawkes) is recorded as \code{shotnoise\_skipped\_in\_prototype} and the
        subset is dropped (\file{grouped\_integral.py:736--757}). The per-diagram
        path accounts for these in \code{delta\_contributions}, which the grouped
        path does not rebuild --- so the wrapper's Tier~2 falls the whole group back
        to per-diagram.
  \item[\code{delta\_contributions} and \code{eval\_per\_diagram\_batch} are
        unsupported.] The grouped return sets them to \code{[]} and \code{None}
        (\file{\_grouped\_phase\_j.py:246--248}). Any caller that needs a
        per-diagram breakdown or the $\delta$-spike accounting must use the
        per-diagram path.
  \item[The group must be homogeneous in compensation.] If
        \code{external\_wick\_compensation} differs across group members the
        evaluator returns \code{status='failed'}
        (\file{grouped\_integral.py:437--448}); the wrapper's Tier~3 then falls
        back to per-diagram. (The wrapper already subdivides by external-legs to
        make this rare.)
  \item[The batch is serial.] As noted, \code{total\_C\_batch} ignores
        \code{parallel}/\code{n\_workers} --- the safe default on the user's
        machine.
\end{description}

\begin{note}
The convention details that must match exactly between the grouped and
per-diagram paths, or boundary contributions would shift: $\Theta(0) = 0$ (a
constraint $\Delta t > 0$ is \emph{strict}, the boundary is outside --- enforced by
the \code{c\_eff <= 0} rejection in both
\code{\_evaluate\_grouped\_m0\_modesum} at \file{grouped\_integral.py:155} and
\code{\_integrate\_grouped\_m1\_modesum} at \file{grouped\_integral.py:217}); and
the overflow guard at $|\gamma_{\text{real}}| > 600$
(\code{\_GROUPED\_EXP\_REAL\_LIMIT}, \file{grouped\_integral.py:135}), a deliberate
margin below the $709$ hard ceiling of \code{cmath.exp} on IEEE doubles, which
returns \code{None} (to fall back) rather than raising.
\end{note}

\section{Where this sits in the pipeline}

To close, the data flow in one picture. The dispatch site is
\code{compute\_cumulants} in \file{pipeline/compute.py}. For each loop order
$\ell$ it gathers the typed diagrams and their prefactors, and if
\code{use\_grouped\_phase\_j} is set it calls the grouped wrapper instead of the
per-diagram path (\file{pipeline/compute.py:757--774}):
\begin{lstlisting}
if use_grouped_phase_j:
    from pipeline._grouped_phase_j import compute_correction_td_grouped
    td_result_ell = compute_correction_td_grouped(
        typed_diagrams = [r['typed_diagram'] for r in records_ell],
        prefactors     = [r['combined_prefactor'] for r in records_ell],
        k              = k,
        propagator_data= propagator_data,
        external_fields= external_fields,
        num_params     = num_params,
        origin_leaf_idx= origin_leaf_idx,
    )
else:
    td_result_ell = compute_correction_td(...)   # the per-diagram path
\end{lstlisting}
Because both paths return the same dict shape, the lines that follow ---
\code{total\_C\_by\_ell[ell] = td\_result\_ell['total\_C']} and the $\tau$-grid
evaluation \code{td\_result\_ell['total\_C\_batch'](tau\_points, ...)}
(\file{pipeline/compute.py:787--794}) --- never learn which path produced the
numbers. The master \code{total\_C} sums across loop orders and $C(\tau) =
\sum_\ell C_\ell(\tau)$ exactly as before. The grouped bookkeeping keys are
available for diagnostics but are not required downstream.

That is the whole subsystem: recognise that a prediagram's typed diagrams share a
domain, sum their integrands once (\eqref{eq:linearity}), pre-sum the residues
into $B_{\bm\alpha}$ (\eqref{eq:Bdef}), integrate the shared exponential-of-linear
in closed form by $m$, divide by the orbit--stabilizer compensation, and fall back
to the per-diagram path the moment any of that cannot be trusted. The payoff is
one JIT compile instead of hundreds and a precision floor at machine $\varepsilon$,
for numbers that match the per-diagram path bit-for-bit wherever the contract
holds.
